% !PRINTED 742
\weberpage{742}

\S.\,202.

\gesperrt{Positive Einheiten.}

Die zuletzt gefundenen Resultate zeigen, dass das System
unabhängiger Einheiten \S.\,201, (7)
\[ \varepsilon'_1, \varepsilon'_2, \ldots, \varepsilon'_{\nu'-1}, \; \mathcal{E}_0, \mathcal{E}_1, \ldots, \mathcal{E}_{\nu'-1} \]
im Körper $H_m$ kein \gesperrt{Fundamentalsystem} bildet. Denn sonst
müssten in den Formeln \S.\,201, (13) die Zahlen
% [sic] -- the source omits a separator between the coefficient families $m_{k,i}$ and $M_{k,i}$; the printed form is preserved.
$m_{k,i} M_{k,i}$ sämmtlich durch 2 theilbar sein, und es wäre
die Determinante $M$, deren Werth wir $= 2^{\nu'}$ gefunden
haben, durch $2^{\nu-1}$ theilbar.

Wir betrachten nun an Stelle des Systems der Gleichungen (18),
\S.\,201 die folgenden Congruenzen:
\weq{1}{
\begin{aligned}
\sum_{1,\,\nu-1}^{i} b_i m_{1,i} &\equiv 1 \\
\sum_{1,\,\nu-1}^{i} b_i m_{2,i} &\equiv 0 \pmod{2}, \\
& \cdots\cdots\cdots\cdots \\
\sum_{1,\,\nu-1}^{i} b_i m_{\nu'-1,i} &\equiv 0
\end{aligned}
}
und beweisen, dass es möglich ist, ihnen durch ganzzahlige
Werthe der $b_i$ zu genügen.

Wäre es nicht möglich, so müsste sich aus den $\nu'-2$
Congruenzen
\[ \sum_{1,\,\nu-1}^{i} b_i m_{k,i} \equiv 0 \pmod{2} \quad k = 2, 3, \ldots, \nu'-1 \]
die letzte
\[ \sum_{1,\,\nu-1}^{i} b_i m_{1,i} \equiv 0 \pmod{2} \]
als \gesperrt{Folgerung} ergeben, und dann würde, wie im vorigen
Paragraphen, weiter folgen:
\[ \sum_{1,\,\nu-1}^{i} b_i M_{k,i} \equiv 0 \pmod{2} \quad k = 0, 1, \ldots, \nu'-1. \]
